% cap1.tex
\chapter{Introducción}
\label{ch:intro}

\textcolor{blue}{En esta sección se debe incluir el fundamento general de la motivación de la investigación desarrollada en la tesis. También se deben incluir la hipótesis de trabajo y los objetivos general y específicos. Máximo 5 páginas. Use el mismo tamaño, tipo de letra e interlineado de este formato.}

\section{Motivación}
En la actualidad, numerosos problemas de toma de decisiones se modelan como problemas de optimización y se resuelven mediante algoritmos computacionales,
que aun cuando pueden buscar llegar a una misma solución, pueden diferir considerablemente en las estrategias algorítmicas utilizadas. 

Por ellos, para un mismo problema existen múltiples alternativas tanto comerciales como propias, por lo que nace la pregunta: 
¿cuál es el mejor algoritmo para un problema específico?

A lo largo de los años, una respuesta simple a esta pregunta ha sido los \glspl{benchmark}, los cuales permiten comparar diferentes métricas de rendimiento sobre un conjunto
de instancias de prueba. Sin embargo, estos \glspl{benchmark} presentan problemas fundamentales, tales como:
\begin{itemize}
   \item Por lo general, solo toman en consideración métricas como rendimiento promedio o en el peor caso (\textit{worst case}), lo cual puede ocultar fortalezas y 
   debilidades específicas de cada algoritmo.
   \item Para que un \gls{benchmark} sea representativo, se requiere de un conjunto de instancias de prueba lo suficientemente grande y diverso, 
   lo cual puede ser costoso y difícil de conseguir.
\end{itemize}

Estos problemas pueden derivar en que los resultados sean sesgados y no representativos de la realidad. Esta observación no es nueva: 
ha sido reportada en múltiples ocasiones y ha llevado a que se busquen nuevas formas de evaluar y comparar algoritmos, en particular existe una 
metodología reciente llamada Análisis del Espacio de Instancias (\textit{Instance Space Analysis}, ISA) desarrollada por Smith-Miles y Muñoz 
\citeyearpar{smith-milesInstanceSpaceAnalysis2023}, 
la cual se basa en el trabajo de Rice \cite{riceAlgorithmSelectionProblem1976}. Esta metodología busca caracterizar y comparar algoritmos a través de un 
espacio de instancias, 
el cual se genera a partir de un conjunto de características de las instancias (\textit{features}) y permite visualizar y analizar el rendimiento de los 
algoritmos en distintas regiones de dicho espacio.

Si bien esta metodología ha sido aplicada a diversos problemas de optimización en el pasado, existe un vacío en la literatura respecto a su aplicación a 
problemas de programación semidefinida (SDP), los cuales son de gran relevancia en la actualidad, puesto que son utilizados para resolver problemas de optimización 
combinatoria, control, aprendizaje automático, entre otros. 
Por otra parte, en la actualidad, no existen muchas fuentes de instancias de prueba para problemas de programación semidefinida, lo cual dificulta la evaluación y 
comparación de algoritmos para este tipo de problemas.

Esto permite plantear la siguiente pregunta de investigación: 
¿Es posible aplicar la metodología ISA a problemas de programación semidefinida?
Y, si es así, ¿es posible llevar un paso más allá esta metodología para guiar la generación de nuevas instancias?


\section{Hipótesis}

\textcolor{blue}{La hipótesis abordada por esta tesis se divide en las siguientes preguntas de investigación:}
\begin{itemize}
   \item PI1:
   \item PI2:
\end{itemize}


\section{Objetivos}
\subsection{Objetivo general}
\subsection{Objetivos específicos}
\begin{enumerate}[]
\item
\item
\item
\end{enumerate}

\section{Estructura de la tesis}
\textcolor{blue}{Esta tesis está organizada en cinco capítulos, cada uno de los cuales contribuye de manera particular al desarrollo de la investigación:}
\begin{enumerate}
\item Capítulo \ref{ch:intro}:
\item Capítulo \ref{ch:litRev}:
\item Capítulo \ref{ch:metodo}:
\item Capítulo \ref{ch:resultados}:
\item Capítulo \ref{ch:conclu}:
\end{enumerate}
